\section{Folding on the original arithmetic quotient}
\label{sec:folding-v059}
Use the notation of QPR1 and of PGR7--PGR13. Here $F$ is the totally
real field, $E/F$ the CM field, $n$ the matrix size, and $q_v$ the
negative index at $v$. Write $S_\infty=N_\infty H_\infty$ for the
real stabilizer, $K=X.K_2$ for the finite level, and $W=W_u$.
Set
\[
  L(g)=\prod_{v\mid\infty}\|g_v\|_{\rm op},\quad
  \nu(x)=\inf_{s\in S_\infty}L(sx),\quad
  B=4n^2\sqrt{2\max_v q_v}.
\]
These are the norm, quotient weight and exponent of the original
$\mathrm{WuDecay}$. Its two premises for the fixed word supply a
\emph{compactly supported} reproducing kernel $\alpha$ and finiteness
of $\int \|W(x)\|\nu(x)^B\,d(X.\nu_Q)(x)$.

% BEGIN FLD1
\paragraph{FLD1. Counting rational translates with the prescribed exponent.}
\label{claim:FLD1}
For all compact $\Omega,C_0\subset G_\infty\times G_f$ there is a
finite constant $A$ such that
\[
 m_\Omega(c(g,1))\le A L(g)^B\qquad(c\in C_0,\ g\in G_\infty).
\]
Here $m_\Omega$ is the supremum over the second argument in (PGR.22).

\begin{proof}
Put $x=c(g,1)$. Fix that second argument $y$. If the counted set
is nonempty choose $\gamma_0$ in it. Right division by $\gamma_0$
injects it into $\Gamma\cap x\Omega\Omega^{-1}x^{-1}$, independently
of $y$. The finite components of the latter elements lie in one fixed
compact set $C_f$, the projection of $C_0$ and $\Omega$ being compact.
The stabilizer $K_0$ of the original finite reference lattice is open
(it contains the open first principal congruence subgroup), so finitely
many translates of $K_0$ cover $C_f$. These matrices have one common
denominator. Explicitly, if $\widehat L$ is the completed reference
lattice and $A_f^0$ the integral coefficient ring, finite generation
and full span give
$d(A_f^0)^n\subset\widehat L\subset d^{-1}(A_f^0)^n$ for one positive
integer $d$: clear the coordinates of the given generators and
of expressions for the coordinate vectors in those generators.
A matrix preserving $\widehat L$ therefore has entries in
$d^{-2}A_f^0$. Clear the finitely many coset representatives as
well: an individual coefficient is a finite sum of pure tensors,
whose number-field and finite adelic factors each have bounded
denominators. For a diagonal rational coefficient $b\in E$, this
implies $D b\in\mathcal O_E$ for a common integer $D$, since it
is integral in every finite completion of $E$ after multiplication
by $D$. The integral adelic tensor ring maps into those local
integer rings. This argument only uses the given congruence topology.

At a real place $v$, use the supplied complex matrix coordinates.
Changing coordinates by the fixed $P_v$, the compact projections
of $C_0,\Omega\Omega^{-1}$ and inversion on $C_0$ only contribute
a constant. In signature coordinates $g_v^{-1}=J_v g_v^*J_v$
with $J_v$ a diagonal isometry, hence
$\|g_v^{-1}\|_{\rm op}=\|g_v\|_{\rm op}\ge1$ ($n>0$).
Each entry of the rational matrix thus has
$|\iota_v(b)|\le C_v\|g_v\|_{\rm op}^2$ with $C_v\ge1$.
Each $v$ selects one member of a complex conjugate pair. The
Minkowski embedding of $D^{-1}\mathcal O_E$ is a full lattice in
$\prod_{v\mid\infty}\mathbb C$. The number of its points in disks
of radii $R_v\ge1$ is bounded by $C\prod_v R_v^2$: translate one
bounded fundamental parallelotope about every such point and compare
its volume with that of the product of disks enlarged by its fixed
coordinate diameters. Taking $n^2$ entries bounds the number of
matrices by $C' L(g)^{4n^2}$. A rational matrix determines both
adelic components, so the same estimate counts group elements.
If some $q_v>0$ then $B\ge4n^2$; if all are zero every $g_v$ is
unitary for the positive form and $L(g)=1$. Both cases give the
claimed exponent, including $B=0$.
\end{proof}
% END FLD1

% BEGIN FLD2
\paragraph{FLD2. The weighted integral on every rational stabilizer domain.}
\label{claim:FLD2}
For every compact $\Omega$, every column $T$, and every original
fundamental domain $F_S$ for the rational stabilizer, one has
\[
             \int_{F_S}m_\Omega(y)|\Phi_u(T)(y)|\,d\mu(y)<\infty
\]
under the original weighted wave hypothesis.

\begin{proof}
Index the finite column classes by $\rho$, write $a_\rho=h_\rho x_u$,
and let $N_\rho$ denote PGR9's integral rational unipotents for the
row lattice $\Lambda_\rho$. We need a compact covering set in the
\emph{real} unipotent group, stronger than existence of a measurable
domain in $G_\infty$. Here is a direct construction. Regard the
Lie algebra of the rational unipotent as a nilpotent $\mathbb Q$-Lie
algebra by restriction of scalars. Its real points are precisely the
product of its points over the infinite places. Exponential and
logarithm are inverse triangular matrix polynomials by PGR8.
Choose a rational basis $X_1,\ldots,X_m$ such that
$[\mathfrak n,X_i]\subset\langle X_1,\ldots,X_{i-1}\rangle$.
Such a basis is obtained successively by choosing a central rational
line in each nonzero nilpotent quotient and taking inverse images.
Scale the basis vectors by positive integers so that both
$\exp(X_i)$ and $\exp(-X_i)$ preserve the given row lattice:
in any $\mathbb Z$-basis of that lattice their finite exponential
polynomials have integral coefficients after clearing the rational
matrix entries and factorials in the nonconstant terms. Thus these
elements belong to $N_\rho$. Reducing a real unipotent by their
integer powers on the left gives a compact set of representatives.
Indeed the ordered exponential coordinates
$\exp(t_1X_1)\cdots\exp(t_mX_m)$ exist uniquely, successively
passing to the last one-dimensional quotient and its kernel.
Left multiplication by $\exp(kX_i)$ adds $k$ to coordinate $i$
and changes only lower coordinates, by the displayed bracket rule.
Reduce coordinates in decreasing order into $[0,1]$. Consequently
$N_\infty=N_\rho C_\rho^N$ for a compact $C_\rho^N\subset N_\infty$.
This includes the zero-dimensional group.

The residual Levi $H_\infty$ fixes both hyperbolic frames and the
negative vector. It acts on their orthogonal complement, a positive
definite Hermitian space at every infinite place (the signature
calculation in QPR5), and is a closed subgroup of a product of
compact unitary groups. PGR2 gives
$S_\infty=N_\rho C_\rho$ with $C_\rho=C_\rho^N H_\infty$ compact.
For any $t\in S_\infty$ write $t^{-1}=n c$ with $n\in N_\rho$,
$c\in C_\rho$. Its unique rational lift $\tilde n$ satisfies
$\tilde n_f^{-1}a_\rho\in a_\rho K$ by PGR9. Therefore, for
every $k\in K$,
\[
    (x,a_\rho k)=\tilde n\,
          (c,\tilde n_f^{-1}a_\rho k)\,(t x,1).
\]
Left rational translation preserves $m_\Omega$ by reindexing.
Apply FLD1 with the compact $C_\rho\times a_\rho K$ and take the
infimum over $t$. This gives $m_\Omega(x,a_\rho k)\le A_\rho\nu(x)^B$.
The inequality holds throughout the finite coset, with no claim that
multiplicity is right $K$-invariant. On that coset PGR6--PGR7 also
give $|\Phi_u(T)(x,a_\rho k)|\le\|J_u T(\rho)\|\|W(x)\|$.

Choose a Borel strict domain $U_\rho$ for $N_\rho$ in $G_\infty$ as
in PGR8. We compare
$\int_{U_\rho}\nu(x)^B\|W(x)\|\,d\mu_\infty(x)$ with the original
quotient integral. The weight $w(x)=\nu(x)^B\|W(x)\|$ is
$S_\infty$-invariant by reindexing the infimum and the unitary wave
covariance; it is Borel since the infimum of continuous functions
is upper semicontinuous, and $\nu\ge1$. Apply the supplied Weil
formula to $\mathbf1_{U_\rho}w$. For fixed $h\in H_\infty$ and
$x\in G_\infty$, the set $\{n\in N_\infty:n h x\in U_\rho\}$
is a Borel fundamental domain for $N_\rho$ in $N_\infty$.
All such domains have the same measure by the countable intersection
partition. It is finite: the compact covering set $C_\rho^N$ above
meets every orbit, and a measurable strict domain can be moved into
it by partitioning according to the first hit in an enumeration
of $N_\rho$. These translated pieces are disjoint.

The measures supplied as $X.\mu_N$ and $X.\mu_H$ restrict to genuine
finite-on-compact Haar measures. To see this even from their printed
weaker description, each is nonzero, sigma-finite, supported on its
closed subgroup, and invariant under its left translations. Both
subgroups are unimodular (compact for $H_\infty$, nilpotent for
$N_\infty$, whose adjoint determinants are one). If $m$ is this
restricted measure and $m_0$ a usual Haar measure, Fubini applied to
$f(y)h(zy)$ under $dm(y)dm_0(z)$ gives
$m(f)m_0(h)=m(h)m_0(f)$, by translation invariance of $m$ and
unimodularity of $m_0$. Choose strictly positive measurable $f$
with finite integral under both, using their sigma-finite exhaustions.
Thus $m$ is a finite positive multiple of $m_0$. In particular the
measure of $H_\infty$ is finite. The original Weil identity now gives
\[
 \int_{U_\rho}w\,d\mu_\infty
   =\operatorname{covol}_{\mu_N}(N_\rho)\,
             \mu_H(H_\infty)\int w\,d(X.\nu_Q)<\infty .
\]
This uses exactly its given normalization and the weighted wave premise.
For arbitrary compact $\Omega$, $m_\Omega$ is Borel by the countable
tuple and compact projection argument of PGR11; that measurability
argument applies also when $1\notin\Omega$. By the PGR10 support
partition, integration of the nonnegative
invariant integrand over $F_S$ equals the sum over the cells
$U_\rho\times a_\rho K$ by left translations. For this step retain
the integral over each whole finite coset, since $m_\Omega$ need
not be right $K$-invariant. Its uniform bound on that coset costs
just the finite factor $\mu_f(K)$. Sum over the finitely many $\rho$.
This proves the claim on every allowed $F_S$.
\end{proof}
% END FLD2

% BEGIN FLD3
\paragraph{FLD3. The original folding implication.}
\label{claim:FLD3}
For every original $D,X,Y,u$, level normalization and
$\mathrm{WuDecay}(X)$ imply $\mathrm{FoldOK}(X,u,Y.hK)$.

\begin{proof}
PGR8 gives a domain for $S_u(F)$ in $G$: the rational subgroup is
discrete and countable and the first-hit construction gives a Borel
strict domain, hence an admissible measure fundamental domain.
For the two uniform quantifiers of the averaging inequality, take
the compact $\Omega_i\ni1$, kernel $\alpha$ and bound $a_i$ supplied
by $\mathrm{WuReprod}$ for this word. Let
$\Omega_0=\Omega_i\times K$ and $a=a_i/\mu_f(K)$; ASG14 gives
$0<\mu_f(K)<\infty$. With the finite argument of the seed held fixed,
it is zero outside its support and on support is a fixed complex
linear evaluation of the archimedean wave, times a unit scalar.
Push that functional through the supplied Bochner reproduction
identity. Taking absolute values bounds the seed by $a_i$ times
the integral of its absolute value over right translations by
$\Omega_i$. Right multiplication of its finite argument by $K$
preserves absolute values, by PGR7, so averaging there rewrites
the bound as
\[
      |\Phi(y)|\le a\int_{\Omega_0}|\Phi(y\omega)|\,d\mu(\omega).
\]
The same $a$ and $\Omega_0$ work for every $T$, and the inequality
persists for every compact $\Omega\supset\Omega_0$ by monotonicity.
Measurability and rational left invariance are PGR7; FLD2 gives the
multiplicity integral for every compact $\Omega$ and every choice
of domain. These are exactly the two clauses of $\mathrm{FoldOK}$.
In particular QPR3 applies at each word under the additional
$L^1$ condition on $X.\mathrm{ph}$, and QPR7--QPR8 apply to a
pair on this same $X,Y$. The compact averaging kernel used here
comes from $\mathrm{WuDecay}$ and need not equal $X.\mathrm{ph}$.
\end{proof}
% END FLD3
